\chapter{Многопоточность}

\section{Основные термины}

\begin{definition}[Процесс]
	\emph{Процесс (process) / задача (task)} -- это цепочка ....
\end{definition}

\begin{remark}
	Процессы выполняются (псевдо) параллельно: каждому выделяются по-очереди квант времени, в течение которого процесс может работать.
\end{remark}

\subsection{Жизенный цикл процесса (общая ситуация)}

Состояние процесса:
\begin{itemize}
	\item новый (процесс только что создан, но еще не встал в очередь);
	\item готовый к выполнению (находится в очереди);
	\item активный -- выполняющийся процесс;
	\item завершенный -- процесс, который завершился, но часть его структур данных еще существует;
	\item блокированный процесс -- процесс, который не может выполняться по причине ожидания какого-то события или ресурса.
\end{itemize}

\begin{definition}[Поток управления]
	\emph{Поток управления (thread)} -- это цепочка выполняющихся команд в рамках какого-то процесса. Поток использует ресурсы процесса.

	Все потоки процесса разделяют ресурсы процессора.
\end{definition}

\begin{remark}
	Процессу выделяется адресное пространство памяти. Адресное пространство процесса может быть больше физической памяти (виртуальная память).
\end{remark}

При переключении процессов необходимо сохранить контекст старого процесса и восстановить контекст нового.

Потоки чередуются, как и процессы. При этом переключение потоков выполняется быстрее, чем у процессов. Также передача данных между потоками быстрее.


\begin{eg}\leavevmode
	\begin{itemize}
		\item интерфейс и сборка проектов в IDE выделяются в отдельные потоки;
		\item фоновое автосохранение;
	\end{itemize}
\end{eg}

\begin{definition}[Атомарная операция]
	\emph{Атомарная операция} -- операция, внутри которой невозможно переключение процессов / потоков.
\end{definition}

Java позволяет создавать потоки и выполнять их параллельно. Поток в Java коде представлен объектом класса \verb|Thread|.

В выполняющейся программе при старте создается главный поток, все остальные создаются во время работы программы. \\

Создание потоков:
\begin{description}
	\item[Способ 1.] Создать объект класса \verb|Thread| специальным конструктором по объекту \verb|Runnable|. Интерфейс \verb|Runnable| описывает действия, коиторые должен выполнять поток (задачу).
	      \begin{lstlisting}[language=Java]
public interface Runnable {
  void run();
}
    \end{lstlisting}
	\item[Способ 2.] Расширить класс \verb|Thread|, переопределив его метод \verb|run()|, создать объект конструктором по умолчанию.

	      После создания объект-потока надо вызвать его метод \verb|start()|.
\end{description}

\newpage

\begin{eg}\leavevmode
	\begin{lstlisting}[language=Java]
public class Demo {
  public static void main(....) {
    S.o.p.("Start");
    Thread thread = new Thread ( () -> {
      for (i = 3; i > 0; i--) {
        S.o.p.("(" + i + ")");
      }
      S.o.p.("Launch");
    }
  );
  thread start();
  S.o.p.("Finish");
}
  \end{lstlisting}
\end{eg}

При создании потоку можно назначить имя.
\[
	\verb|Thread.currentThread| \ \text{-- текущий поток}.
\]

\begin{lstlisting}[language=Java]
public static main(...) {
  S.o.p.("Start");
  Runnable task = () -> {
    String name = Thread.currentThread.getName();
    for (int i = 3; i > 0; i--) {
      S.o.p.(name + "(" + i + ")");
    }
    S.o.p.(name + "launch!");
  };
}
\end{lstlisting}
